En el intervalo \(I=(-1,1)\), considere
\[
u(x)=x^3,\qquad v(x)=|x|^3.
\]
Compruebe que ambas funciones son de clase \(C^2\), calcule su wronskiano en todo \(I\) y decida si son linealmente independientes. ¿Podrían ser ambas soluciones de una misma ecuación
\[
y''+p(x)y'+q(x)y=0
\]
con \(p,q\) continuas en \(I\)? Justifique la respuesta. Explique después por qué, para dos soluciones de una misma ecuación con esas hipótesis, que el wronskiano sea cero en un solo punto sí implica dependencia lineal.
